\color{unidarkgreen}{\subsection{Орган выявления БНТ ВВ}}
\color{black}

\begin{enumerate}[label=\arabic{section}.\arabic{subsection}.\arabic*,  align=left, labelsep=4pt, leftmargin=0pt, itemindent=57pt]

\item
Для обеспечения несрабатывания «МТЗ ВВ», «ТЗНП ВВ», «ТЗНП НВ», «ТЗОП ВВ» при броске тока
намагничивания используется блокировка по содержанию 2 гармоники в токе. Блокировка основана на
контроле отношения действующего значения второй гармонике к первой тока фаз. Имеется возможность
регулировать порог срабатывания блокировки, с помощью уставки \textbf{«I2г/I1г БНТ ВВ»}. Блокировка по 2
гармонике выводится из действия при величине тока ниже 0,04Iном, при срабатывании \textbf{«Iср БНТ»}, а так
же при выводе программной накладкой \textbf{«БНТ ВВ»}.  

\begin{figure}[H]
\centering
\includegraphics[width=0.7\textwidth,height=0.7\textheight,keepaspectratio]{\fbpath/171.0391 БНТ/img/BNT.pdf}
\captionof{figure}{Функциональный блок БНТ}\label{pict_bnt:BNT}
\end{figure}

\small
\begin{longtable}{|>{\centering\arraybackslash}m{5.3cm}|>{\centering\arraybackslash}m{4.2cm}|>{\centering\arraybackslash}m{1cm}|>{\centering\arraybackslash}m{3.3cm}|>{\centering\arraybackslash}m{1.8cm}|}
% Первая страница: полный заголовок
\caption{Перечень уставочных параметров БНТ\hfill\vspace{-0.5\baselineskip}}\label{bnt:tbl1}\\
\hline
\rowcolor{gray!30}
Наименование уставки & Диапазон & Шаг & Значение по умолчанию & Единица измерения \\
\hline
\endfirsthead

% Промежуточные страницы: «Продолжение таблицы»
\caption*{\hspace{3pt}\emph{Продолжение таблицы \ref{bnt:tbl1}}\hfill\ }\vspace{-0.5\baselineskip} \\
\hline
\rowcolor{gray!30}
Наименование уставки & Диапазон & Шаг & Значение по умолчанию & Единица измерения \\
\hline
\endhead

%===m>БНТ|-
%===m
\end{longtable}
\normalsize

\end{enumerate}